% SSSOM/T and Uberon presentation
% © 2024 Damien Goutte-Gattat
% 
% This work is released under the Creative Commons Attribution 4.0
% International License. To view a copy of this license, visit
% https://creativecommons.org/licenses/by/4.0/.
\usepackage{fontspec}
\usepackage{polyglossia}
\usepackage{tikz}

\setmainlanguage{english}
\usetheme{cambridge}

\graphicspath{{imgs/}{svgs/}}

% TikZ libraries and styles
\usetikzlibrary{positioning,shapes.geometric,graphs,graphdrawing,quotes,fit,overlay-beamer-styles}
\usegdlibrary{layered,trees}
\colorlet{GenericItem}{yellow!20!white}
\colorlet{GenericHighlight}{red!20!white}
\tikzset{Term/.style={draw,fill=yellow!20!white,font=\tiny\bfseries}}
\tikzset{Relation/.style={draw=blue!40!white,thick,arrows=<-}}
\tikzset{RelationLabel/.style={shape=circle,fill=white,draw=black,thin,font=\tiny\bfseries,inner sep=1pt}}
\tikzset{File/.style={draw,fill=yellow!20!white,font=\tiny\ttfamily}}
\tikzset{Entity/.style={draw,fill=yellow!20!white,font=\scriptsize}}
\tikzset{Process/.style={draw,shape=ellipse,fill=blue!20!white,font=\scriptsize}}

\title{(Re-)bridging the anatomy ontologies with SSSOM}

\author{Damien~Goutte-Gattat\inst{1}}
\institute{\inst{1}FlyBase group, Department of Physiology, Development and Neuroscience}
\date{International Conference on Biomedical Ontologies 2024, July 2024}

\hypersetup{pdfpagemode=UseNone,
  pdftitle={},
  pdfauthor={Damien Goutte-Gattat}}
\AtBeginSection{
  \begin{frame}{Outline}{}
    \tableofcontents[currentsection]
  \end{frame}
}

\begin{document}

\begin{frame}
  \titlepage
\end{frame}

\section{Introduction}

\begin{frame}
  \frametitle{An approach to a multi-species anatomy ontology}

  \begin{columns}
    \begin{column}{.5\textwidth}
      \begin{tikzpicture}[nodes={shape=circle,draw,fill=GenericItem}]
        \node (core) [alt=<6| handout:0>{fill=GenericHighlight}{}] {Uberon};
        \node (fbbt) [alt=<{1,7| handout:0}>{fill=GenericHighlight}{}] at (150:2cm) {FBbt};
        \node (zfa)  [alt=<{2,7| handout:0}>{fill=GenericHighlight}{}] at (90:2cm)  {ZFA};
        \node (xao)  [alt=<{3,7| handout:0}>{fill=GenericHighlight}{}] at (30:2cm)  {XAO};
        \node (ma)   [alt=<{4,7| handout:0}>{fill=GenericHighlight}{}] at (210:2cm) {MA};
        \node (wbbt) [alt=<{5,7| handout:0}>{fill=GenericHighlight}{}] at (270:2cm) {WBbt};
        \node (tx1) at  (310:2cm) {...};
        \node (tx2) at  (330:2cm) {...};
        \node (tx3) at  (350:2cm) {...};
        \draw (core) -- (fbbt) (core) -- (zfa)  (core) -- (xao)
              (core) -- (ma)   (core) -- (wbbt) (core) -- (tx1)
              (core) -- (tx2)  (core) -- (tx3);
      \end{tikzpicture}
    \end{column}
    \begin{column}{.45\textwidth}
      Many species-specific anatomy ontologies available:
      \smallskip
      \begin{overlayarea}{\textwidth}{1cm}\color{red}
        \only<1>{\emph{Drosophila} Anatomy Ontology}
        \only<2| handout:0>{Zebrafish Anatomy Ontology}
        \only<3| handout:0>{Xenopus Anatomy Ontology}
        \only<4| handout:0>{Mouse Anatomy Ontology}
        \only<5| handout:0>{\emph{C. elegans} Anatomy Ontology}
      \end{overlayarea}

      Strategy for a multi-species ontology:
      \begin{itemize}
        \item Provide a \textcolor<6>{red}{taxon-neutral core framework}
        \item Integrate \textcolor<7>{red}{existing taxon-specific ontologies}
      \end{itemize}
    \end{column}
  \end{columns}

  \note<1>{Anatomy ontologies are used in many aspects of biocuration, for
  example to capture the organ of origin of a biological sample, the tissue
  where a gene is expressed, or the anatomical location where a disease
  manifests itself.

  Many species-specific curation projects have, therefore, their own ontology
  to describe the anatomy of their favourite organism. For example, at FlyBase
  we have our own \emph{Drosophila} Anatomy Ontology (FBbt), ZFIN has its
  Zebrafish Anatomy Ontology (ZFA), XenBase has its Xenopus Anatomy Ontology
  (XAO), and so on.

  But to enable \emph{cross-species} curation and analysis of data (especially
  high-throughput data, such as scRNAseq), we need ideally a single,
  multi-species ontology that covers all our model organisms. For the past 12
  years or so, this has been the goal of Uberon, the Uber Anatomy Ontology.

  Uberon's strategy to achieve that goal is to provide a taxon-neutral core
  framework, describing common anatomical structures in broadly generic terms,
  and to plug the species-specific ontologies into that core.}

\end{frame}

\begin{frame}
  \frametitle{Example: integrating FBbt into Uberon}

  \only<1| handout:1>{
    \begin{columns}[t]
      \begin{column}{.45\textwidth}
        Uberon hierarchy
        \begin{tikzpicture}
          \graph[layered layout,orient=down,nodes={Term},edge quotes mid,edges={Relation,nodes={RelationLabel}}] {
            anatomical entity <-
            anatomical structure <-
            circulatory organ <- {
              primary circulatory organ <- { heart, dorsal vessel heart[fill=GenericHighlight] },
              accessory circulatory organ <- lymph heart
            };
          };
        \end{tikzpicture}
      \end{column}
      \begin{column}{.45\textwidth}
        FBbt hierarchy
        \begin{tikzpicture}
          \graph[layered layout,orient=down,nodes={Term},edge quotes mid,edges={Relation,nodes={RelationLabel}}] {
            anatomical entity <-
            anatomical structure <-
            heart[fill=GenericHighlight] <-["P"] {
              aortic funnel, heart chamber, dorsal diaphragm
            };
          };
        \end{tikzpicture}

        FBbt's ``heart'' is the fly-specific equivalent of Uberon's ``dorsal vessel heart''
      \end{column}
    \end{columns}
  }
  \only<2| handout:2>{
    \begin{columns}
      \begin{column}{.3\textwidth}
        Integrated hierarchy:

        \smallskip
        FBbt's ``heart'' classified under Uberon's ``dorsal vessel heart''
      \end{column}
      \begin{column}{.6\textwidth}
        \begin{tikzpicture}
          \graph[layered layout,orient=down,nodes={Term},edge quotes mid,edges={Relation,nodes={RelationLabel}}] {
            anatomical entity <-
            anatomical structure <-
            circulatory organ <- {
              primary circulatory organ <- {
                heart,
                dorsal vessel heart[fill=GenericHighlight]
                    <- "heart (drosophila)"[fill=GenericHighlight] <-["P"] {
                  aortic funnel, heart chamber, dorsal diaphragm
                }
              },
              accessory circulatory organ <- lymph heart
            };
          };
        \end{tikzpicture}
      \end{column}
    \end{columns}
  }
  
  \note<1>{To illustrate this strategy, here is (on the left) a simplified
  excerpt of Uberon's hierarchy around the concept of ``heart'', which in
  Uberon parlance is called, more generically, the ``primary circulatory
  organ'' (to make room for the \emph{accessory} circulatory organs that can
  be found in some species).

  Below ``primary circulatory organ'', Uberon defines both ``heart'', which is
  intended to refer to the \emph{vertebrate} heart, and ``dorsal vessel
  heart'', which refers to the heart of arthropods such as fruit flies. Uberon
  itself does not describe the arthropod heart in any more details -- there
  are no terms under ``dorsal vessel heart''. That term is in fact merely
  intended as a ``branching point'', where more specific terms from
  arthropod-specific ontologies can be rattached.

  In the \emph{Drosophila} Anatomy Ontology, we have a single term ``heart''
  to refer to the dorsal vessel heart (we don't need to call it ``dorsal
  vessel heart'' since it is the only heart that exists in flies), and below
  that term, we have a handful of terms that describe more precisely the
  different parts that make up the fly heart.}

  \note<2>{When FBbt is integrated into Uberon, its ``heart'' term is
  automatically treated as a subclass of the more generic ``dorsal vessel
  heart'', thereby creating a unified hierarchy where the
  \emph{Drosophila}-specific terms are appropriately rattached to Uberon's
  taxon-neutral core.

  Using this strategy, Uberon is able to accurately and precisely describe the
  anatomy of the most used model organisms, by benefitting from all the
  existing species-specific ontologies without having to duplicate their
  efforts.

  This talk is about how this integration process works, and how we recently
  overhauled it to make it easier to maintain, more robust, and FAIRer.}

\end{frame}

\begin{frame}
  \frametitle{Integration requires cross-ontology mappings}

  \begin{block}{}
    Integrating the taxon-specific ontologies with the Uberon core framework
    requires to know where taxon-specific terms should be ``plugged'' in the
    core framework.

    E.g., we need to know that FBbt's \emph{heart} must be attached to Uberon's
    \emph{dorsal vessel heart}.
  \end{block}

  \note{The key component of the integration process is the knowledge of which
  species-specific terms, from each of the species-specific ontologies such as
  FBbt, should be ``attached'' to which taxon-neutral term in Uberon's core
  framework. For example, knowing that FBbt's ``heart'' is the fly-specific
  term for what Uberon calls the ``dorsal vessel heart''.

  This is the role of cross-ontology \emph{mappings}.}

\end{frame}

\section{Limits of the previous system}

\begin{frame}
  \frametitle{Overview of the previous integration system}

  \begin{columns}
    \begin{column}{.46\textwidth}
      \begin{tikzpicture}[node distance=.3cm]
        \node [Entity,text width=1.3cm,alt=<1>{fill=GenericHighlight}{}] (uberon) {Uberon\\(incl. xrefs)};
        \node [Process,below right=of uberon,alt=<3>{fill=GenericHighlight}{}] (script) {Custom script};
        \node [Entity,above right=of script,align=left] (ssao) {Taxon-specific\\
                                                                ontologies};
        \node [Entity,below=of script,alt=<2>{fill=GenericHighlight}{}] (cob) {Bridge files};
        \node [Process,below=of cob] (merge) {OWL merge};
        \node [Entity,below=of merge] (cm) {Multi-species ontology};
        \draw [->] (uberon.east) -| (script.north) (script.south) -- (cob.north);
        \draw [->] (cob) -- (merge) (merge) -- (cm);
        \draw [->] (uberon.south) |- (merge.west);
        \draw [->] (ssao.south) |- (merge.east);
      \end{tikzpicture}
    \end{column}
    \begin{column}{.51\textwidth}
      \begin{block}{Mappings maintained as \textcolor<1>{red}{xrefs}}\scriptsize\ttfamily
[Term]

id: UBERON:0015230

name: dorsal vessel heart

is\_a: UBERON:0007100 ! primary circulatory organ

\textcolor<1>{red}{xref: FBbt:00003154}
      \end{block}
  
      \begin{block}{Xrefs transformed to \textcolor<2>{red}{bridge files}}\scriptsize\ttfamily
[Term]

id: FBbt:00003154

intersection\_of: UBERON:0015230

intersection\_of: part\_of NCBITaxon:7227
      \end{block}
    \end{column}
  \end{columns}

  \begin{itemize}
    \item Mappings not available to third parties
    \item \textcolor<3>{red}{Custom pipeline} highly Uberon-specific, inflexible and not reusable
  \end{itemize}

  \note<1>{Mappings are represented and maintained as \emph{cross-references}
  that are carried by the taxon-neutral terms in Uberon and that point to the
  taxon-specific terms in other ontologies -- as shown here, where Uberon's
  ``dorsal vessel heart'' term carries a ``xref'' to FBbt:00003154, which is
  FBbt's term for the fly heart.

  When building Uberon, those cross-references are extracted by a custom
  script that turns them into \emph{bridge files}, which contains OWL
  materialisations of the mappings, in the form of equivalence axioms (for
  example here, an axiom stating the FBbt's ``heart'' is equivalent to an
  Uberon's ``dorsal vessel heart'' that is part of the \emph{Drosophila
  melanogaster} taxon).

  This system has severe limitations, one of the most important being that the
  mappings are directly embedded within the ontology itself and are not
  independently available. In FAIR terms, this makes the mappings much more
  difficult to find, access and reuse than we would like.

  Furthermore, the custom script that generates the bridge files from the
  cross-references is highly specific to Uberon; its behaviour are completely
  hardcoded to fit Uberon's use case, so the integration logic is difficult to
  modify and almost impossible to reuse by other projects.}

\end{frame}

\begin{frame}
  \frametitle{Goals for a new approach}

  \begin{enumerate}
    \item Cross-ontology mappings (COM) as first-class entities
      \begin{itemize}
        \item available independently of the ontology itself
      \end{itemize}
    \item Represent mappings using a standard format intended for that purpose
      \begin{itemize}
        \item allowing for mapping-specific metadata
        \item allowing for meaningful mapping predicates
      \end{itemize}
    \item Make it easy to derive OWL bridge files from mappings
      \begin{itemize}
        \item without requiring Uberon-specific tooling
      \end{itemize}
  \end{enumerate}

  \note{In response to the limits of the previous system, we set out to design
  a new system, with the following aims:

  First, we want to treat mappings as ``first-class entities'', on the same
  level as the terms of the ontology. This notably means that they must be
  publicly available on their own, without requiring people interested in them
  to forcefully extract them from the ontology.

  Second, we want to store the mappings using a format that is both standard
  and specifically intended for representing mappings. That format should make
  it possible to store mapping-specific metadata, and to give precise meaning
  to a mapping (not just saying that two terms are mapped, but: who asserted
  that those two terms should be mapped? When? On what basis? And what does
  the mapping actually means?).

  Lastly, we want a bridge generation process that is much simpler than the
  current one, and that is not tied to Uberon in any way, so that it can be
  reused by other projects that have a similar need for integrating several
  ontologies together.}

\end{frame}

\section{Implementation}

\begin{frame}
  \frametitle{The SSSOM standard}

  ``Simple Standard for Sharing Ontological Mappings''
  \begin{itemize}
    \item A data model to represent semantic mappings and associated metadata
    \item A simple TSV-based format to store and exchange them
  \end{itemize}

  \begin{block}{SSSOM/TSV format}\tiny\ttfamily
    \begin{tabular}{lllll}
      \hline
      subject\_id   & predicate\_id                 & object\_id      & mapping\_justification        & mapping\_date\\
      \hline
      FBbt:00000001 & semapv:crossSpeciesExactMatch & UBERON:0000468  & semapv:ManualMappingCuration  & 2024-06-01\\
      FBbt:00000002 & semapv:crossSpeciesExactMatch & UBERON:6000002  & semapv:ManualMappingCuration  & 2024-06-01\\
      ...           & ...                           & ...             & ...                           & ...\\
      \textcolor<2>{red}{FBbt:00003154} & \textcolor<2>{red}{semapv:crossSpeciesExactMatch} & \textcolor<2>{red}{UBERON:0015230}  & \textcolor<3>{red}{semapv:ManualMappingCuration}  & \textcolor<3>{red}{2024-06-01}\\
      ...           & ...                           & ...             & ...                           & ...\\
      \hline
    \end{tabular}
  \end{block}

  \begin{overlayarea}{\textwidth}{1cm}
    \only<2>{Mapping proper: subject -- predicate -- object triple}
    \only<3>{Mapping metadata}
  \end{overlayarea}

  \note<1>{Over the past three years, independently of Uberon, a new standard
  has been devised to work with semantic mappings: the Simple Standard for
  Sharing Ontological Mappings.

  The SSSOM standard defines a common data model to represent and manipulate
  mappings in programming languages. The model notably allows to capture a
  rich set of metadata about each mapping.

  The standard also defines several serialisation formats to store and
  exchange mappings, the most important of which being the ``SSSOM/TSV''
  format, intended to make reusing mappings across resources as easy as
  possible.

  You see here an very basic example of what a mapping set in SSSOM/TSV format
  looks like: one mapping per line, each mapping consisting, at its core, of a
  triple made of a subject, an object, and a mapping predicate. That core is
  associated with several metadata.}

\end{frame}

\begin{frame}
  \frametitle{New mapping predicates for cross-species mappings}

  Additions to the \emph{Semantic Mapping Vocabulary} (SEMAPV)
  \begin{center}
    \begin{tikzpicture}
      \graph[tree layout,nodes={Term},edges={Relation}] {
        "skos:mappingRelation" <- {
          "semapv:isomorphicMatch" <- {
            "skos:exactMatch",
            "semapv:crossSpeciesExactMatch" [fill=GenericHighlight]
          },
          "semapv:nonIsomorphicMatch" <- {
            "skos:broadMatch" <- "semapv:crossSpeciesBroadMatch" [fill=GenericHighlight],
            "skos:narrowMatch" <- "semapv:crossSpeciesNarrowMatch" [fill=GenericHighlight]
          }
        };
      };
    \end{tikzpicture}
  \end{center}

  \emph{crossSpeciesExactMatch}: ``a match where the subject is considered analogous
  to the object in a different taxonomic grouping''

  \note{One of the benefits of the SSSOM standard is that it allows to specify
  an explicit mapping predicate, so as to give a precise meaning to each
  mapping. It is then possible to distinguish, for example between mappings of
  terms that represent exactly the same concept in two different vocabularies,
  and mappings of terms that represent \emph{similar}, but not exactly
  identical, concepts.

  The authors of the SSSOM standard have compiled a dedicated vocabulary, the
  \emph{Semantic Mapping Vocabulary}, to host, among other things, various
  mapping predicates that can be used to qualify mappings. In that vocabulary,
  as part of this project, we have added a handful of predicates to
  specifically represent \emph{cross-species mappings}.

  The most important being \emph{crossSpeciesExactMatch}, which is intended
  for mappings between terms that represent the same concept, modulo a
  taxonomic restriction. We can use that predicate to explicitly say that the
  FBbt term for the fly heart represents exactly the same concept as the
  Uberon term for dorsal vessel heart, \emph{except} for the fact that the
  former is specific to flies.}

\end{frame}

\begin{frame}[t]
  \frametitle{From SSSOM sets to bridge files: SSSOM/T-OWL}

  A domain-specific language to transform SSSOM mappings into OWL axioms

  \begin{block}{Model of a mapping processing rule}
    \begin{center}
      \begin{tikzpicture}
        \node [text width=5cm,fill=GenericItem,draw] (filter) at (0,0)
              {\textbf{Filter}\\
               Mapping $\rightarrow$ boolean\\
               Decide whether the rule applies to a mapping};
        \node [text width=5cm,fill=GenericItem,draw] (preprocessor) at (6,1)
              {\textbf{Preprocessor}\\
               Mapping $\rightarrow$ mapping\\
               Modify a mapping on the fly};
        \node [text width=5cm,fill=GenericItem,draw] (generator) at (6,-1)
              {\textbf{Generator}\\
               Mapping $\rightarrow$ OWL axiom\\
               Derive an axiom from a mapping};
        \draw[->] (filter) -- (preprocessor);
        \draw[->] (filter) -- (generator);
      \end{tikzpicture}
    \end{center}
  \end{block}

  \note{Now that we can store mappings in a dedicated format, along with their
  metadata and with meaningful mapping predicates, we need a way to use those
  mappings when we want to integrate several ontologies together.  And, as per
  our third goal, we want to do that in a way that is, as much as possible,
  not specific to Uberon, but reusable by other projects as well.

  For that, we have designed, and we propose, a small domain-specific language
  specifically intended to process SSSOM mappings and turn them into arbitrary
  OWL axioms. We call that language \emph{SSSOM/T-OWL}, or ``SSSOM Transform
  to OWL''.

  The core concept of SSSOM/T-OWL is the mapping processing rule, which is
  made of two elements.

  First, a filter that takes a mapping and returns a boolean value, indicating
  whether the rule should apply to that mapping or not; the filter can examine
  all the metadata of a mapping to make its decision.

  Second, either a preprocessor that takes a mapping and returns a modified
  version of that mapping (allowing to modify a mapping on the fly), or a
  generator that takes a mapping and returns an OWL axiom.}
\end{frame}

\begin{frame}
  \frametitle{From SSSOM sets to bridge files: SSSOM/T-OWL}

  Examples of SSSOM/T-OWL rules

  \begin{overlayarea}{\textwidth}{3cm}
    \only<1-2| handout:1>{
      \begin{block}{Inverting a mapping}\ttfamily\scriptsize
        \textcolor<1>{red}{subject==UBERON:*} -> \textcolor<2>{red}{invert()};
      \end{block}
    }
    \only<3-4| handout:2>{
      \begin{block}{Generating a simple axiom}\ttfamily\scriptsize
        \textcolor<3>{red}{predicate==skos:broadMatch} -> \textcolor<4>{red}{create\_axiom("\%subject\_id SubClassOf: \%object\_id")};
      \end{block}
    }
    \only<5-10| handout:3>{
      \begin{block}{Generating a more complex axiom}\ttfamily\scriptsize
        \textcolor<5>{red}{subject==FBbt:*}\\
        \hspace{1cm} \&\& \textcolor<6>{red}{predicate==semapv:crossSpeciesExactMatch}\\
        \hspace{1cm} \&\&
        (\textcolor<7>{red}{mapping\_justification==semapv:ManualMappingCuration}
        \textcolor<8>{red}{||} \textcolor<9>{red}{confidence>=0.8})\\
        ->  \textcolor<10>{red}{create\_axiom("\%subject\_id EquivalentTo: \%object\_id and\\
        \hspace{2cm} (BFO:0000050 some NCBITaxon:7227"))};
      \end{block}
    }
  \end{overlayarea}

  \begin{overlayarea}{\textwidth}{2.5cm}
    \only<1-2| handout:1>{\textcolor<1>{red}{Select mappings with a UBERON subject}

    and \textcolor<2>{red}{flip them around the predicate}}
    \only<3-4| handout:2>{\textcolor<3>{red}{Select mappings with a \emph{skos:broadMatch} predicate}

    and \textcolor<4>{red}{create a \emph{SubClassOf} axiom}}
    \only<5-10| handout:3>{Select mappings
      \begin{itemize}
        \item \textcolor<5>{red}{with a FBbt subject}
        \item \textcolor<6>{red}{with a \emph{semapv:crossSpeciesExactMatch} predicate}
        \item that have been \textcolor<7>{red}{manually curated} \textcolor<8>{red}{or} have a \textcolor<9>{red}{confidence higher than 80\%}
      \end{itemize}

      and \textcolor<10>{red}{create an equivalence axiom with a taxonomic restriction}}
  \end{overlayarea}

  \note<1>{To illustrate, here are some examples of SSSOM/T-OWL rules.

  First, a simple rule with a simple filter that selects all mappings where
  the subject is a Uberon term, and apply an inversion preprocessor to them.
  Inversion flips a mapping around the predicate, so that the subject becomes
  the object and the other way round. So after applying this rule, all
  mappings with a Uberon subject in the set will become mappings with a Uberon
  object.

  Second, a rule with an equally simple filter that selects all mappings that
  use a \emph{skos:broadMatch} predicate, and generate from them a SubClassOf
  axiom that states that the subject term is a subclass of the object term
  (which represents a broader concept).

  Third, a more complex rule, with a compound filter that selects mappings
  that: (1) have a FBbt subject, (2) use a \emph{crossSpeciesExactMatch}
  predicate, (3) have either been manually curated, or have a confidence level
  higher than 80\%. When such a mapping is found, the rule creates an
  equivalence axiom stating that the subject (the FBbt term) is equivalent to
  the object that is part of the \emph{Drosophila melanogaster} taxon.}

\end{frame}

\begin{frame}
  \frametitle{New SSSOM-based integration pipeline}

  \begin{columns}
    \begin{column}{.45\textwidth}
      \begin{center}
        \begin{tikzpicture}[node distance=.5cm]
          \node [Entity] (uberon) {Uberon};
          \node [Entity,right=of uberon] (com) {Mapping sets};
          \node [Entity,right=of com,align=left] (ssao) {Taxon-specific\\
                                                         ontologies};
          \node [Process,below=of com,alt=<2>{fill=GenericHighlight}{}] (sssomt) {SSSOM/T-OWL};
          \node [Entity,below=of sssomt] (cob) {Bridge files};
          \node [Process,below=of cob,alt=<3>{fill=GenericHighlight}{}] (merge) {OWL merge};
          \node [Entity,below=of merge] (cm) {Multi-species ontology};
          \draw [->] (com) -- (sssomt) (sssomt) -- (cob);
          \draw [->] (cob) -- (merge) (merge) -- (cm);
          \draw [->] (uberon.south) |- (merge.west);
          \draw [->] (ssao.south) |- (merge.east);
        \end{tikzpicture}
      \end{center}
    \end{column}
    \begin{column}{.45\textwidth}
      \begin{enumerate}
        \item Generating the bridge files through a \textcolor<2>{red}{SSSOM/T-OWL ruleset}\\
              (\texttt{robot sssom:inject})
        \item \textcolor<3>{red}{Merging} Uberon, the bridges, and the taxon-specific ontologies\\
              (\texttt{robot merge})
      \end{enumerate}
    \end{column}
  \end{columns}

  \note<1>{This leads us to our new, SSSOM-based integration pipeline for Uberon.
  That pipeline has three distinct inputs: (1) Uberon itself, (2) the mappings
  between Uberon and the taxon-specific ontologies (which are no longer
  embedded within the ontology itself), and (3) the taxon-specific ontologies.

  The bridge files are generated by a ROBOT plugin that implements the
  SSSOM/T-OWL language, and a custom SSSOM/T-OWL ruleset that contains the
  Uberon-specific logic, in the form of rules similar to the one seen in the
  last example.

  After this step, we have only OWL files, which can all be merged using a
  standard ROBOT merge operation, producing the multi-species ontology that we
  want.

  Overall, that pipeline is both more robust, as most of the code is
  consolidated in ROBOT, and more flexible, as the Uberon-specific logic is
  easily modifiable in the SSSOM/T-OWL ruleset.}

\end{frame}

\section{Conclusion}

\begin{frame}
  \frametitle{Benefits of the new approach}

  \begin{block}{Cross-ontology mappings as first-class artefacts}
    \begin{itemize}
      \item With their own metadata
      \item Can be maintained separately
      \item Available as a separate release product in a standard format
    \end{itemize}
  \end{block}

  \begin{block}{Generic framework for SSSOM-based integration of multiple ontologies}
    \begin{itemize}
      \item SSSOM/T-OWL is completely independent of Uberon
      \item Only the ruleset is specific
    \end{itemize}
  \end{block}

  \note{Beyond the more robust and more flexible pipeline, the new approach
  benefits for downstream users of Uberon as well.

  The mappings between Uberon and the taxon-specific ontologies are now
  first-class artefacts. They can have their own metadata, allowing for
  traceability and accountability of each individual mapping statement. They
  can be maintained separately from the Uberon ontology itself, for example to
  delegate them to the taxon-specific ontologies themselves (something we have
  already started doing for the mappings with FBbt, which are now maintained
  by FBbt rather than by Uberon).  They are available as a distinct release
  product, which is findable on its own, in a format that is still relatively
  new but that we hope will become the standard format for semantic mappings.

  And for everyone else, we now have a generic framework to integrate multiple
  ontologies with SSSOM-maintained mappings. The SSSOM/T-OWL language, and its
  implementation in the SSSOM ROBOT plugin, has no ties whatsoever to Uberon
  and can be reused by any project. Only the SSSOM/T-OWL rules that we use are
  tailored to the needs of Uberon.}

\end{frame}

\begin{frame}
  \frametitle{About}

  \begin{block}{Acknowledgements}
    \begin{itemize}
      \item Grant BB/T014008 from BBSRC (UK) and NSF/BIO (US)
      \item Mapping Commons project for the SSSOM standard
    \end{itemize}
  \end{block}

  \begin{block}{\includegraphics[scale=.26]{ccby}\hspace{.2em}%
    2024 Damien Goutte-Gattat \texttt{<dpg44@cam.ac.uk>}}
    This presentation is released under the Creative Commons Attribution
    4.0 International License (CC BY 4.0).
  \end{block}

  \begin{block}{Revision}\ttfamily\scriptsize
    \input{revision.tex}
  \end{block}
\end{frame}

\appendix
\begin{frame}[noframenumbering]
  \frametitle{Questions?}
\end{frame}

\begin{frame}[noframenumbering]
  \frametitle{Implementation in ROBOT commands}

  \begin{block}{sssom:xref-extract command}
    Convert existing mappings represented as cross-references into a SSSOM set

    \medskip
    {\ttfamily\scriptsize
robot sssom:xref-extract -i input.owl --mapping-file output.sssom.tsv
    }
  \end{block}

  \begin{block}{sssom:inject command}
    Apply SSSOM/T-OWL rules to a mapping set and inject the resulting axioms into an ontology

    \medskip
    {\ttfamily\scriptsize
robot sssom:inject -i input.owl --sssom mappings.sssom.tsv --ruleset bridge.rules
    }
  \end{block}

  \note{To actually implement our new approach, we first, as a preliminary
  step, added plugin support to ROBOT, the main tool we use to manipulate OWL
  and OBO ontologies. This is so that we can quickly add new commands to
  ROBOT. We then developed a SSSOM plugin for ROBOT, which provides two
  distinct commands.

  The first is a command to facilitate the transition to SSSOM-maintained
  mappings by automatically extracting mappings stored as cross-references
  within an ontology into a mapping set in the SSSOM/TSV format.

  More importantly, the second is a command that implements the SSSOM/T-OWL
  domain-specific language. That command takes a mapping set and a set of
  SSSOM/T-OWL rules as input, and will inject into the current ontology the
  axioms produced by applying the SSSOM/T-OWL rules to the mapping set.}

\end{frame}

\begin{frame}[fragile,noframenumbering]
  \frametitle{Example: FBbt/Uberon/CL integration test}

  A single ROBOT command to merge FBbt with Uberon, CL, and the mappings-derived
  bridging axioms, and check for possible inconsistencies that would reveal
  incorrect mappings:

  \begin{block}{}
\begin{verbatim}
robot merge -i fbbt.owl -i uberon.owl -i cl.owl \
      sssom:inject --sssom mappings.sssom.tsv   \
                   --ruleset bridging.rules     \
      reason -r ELK
\end{verbatim}
  \end{block}

  \note{Lastly, here is an example to illustrate some of the benefits of the
  approach for third-party users. As I've mentioned, at FlyBase we are now
  maintaining ourselves the mappings between FBbt and Uberon.

  Because FBbt and Uberon are developed separately, there is always a risk
  that, whenever something changes in either FBbt, Uberon, or the
  FBbt-to-Uberon mappings, an inter-ontology inconsistency is introduced.
  Testing to detect and prevent such inconsistencies is difficult and,
  historically, has never actually been done on a regular basis.

  But with the SSSOM approach, such testing can be done with a single ROBOT
  command in three steps, as shown here: (1) we merge the current versions of
  FBbt, Uberon, and CL; (2) we derive OWL axioms from the FBbt-to-Uberon/CL
  mappings and inject them into the merged ontology; (3) we reason over the
  merged ontology to detect inconsistencies. This takes but a couple of
  minutes, and that test is now part of our test suite in FBbt. Thanks to this
  approach, over the past two years we have rooted out all existing
  inconsistencies between FBbt and Uberon.}

\end{frame}

\end{document}
